% Source: https://artofproblemsolving.com/community/c3h2391932_difficult_problem__average_marks_problem
% Stable identity: aops-topic-2391932-q01
An average times the number of students is the total of their marks. The whole class has total $ax$, the top group $by$ and the bottom group $cz$. The remaining $x-y-z$ students therefore have total $ax-by-cz$, and their average is
\[\boxed{\dfrac{ax-by-cz}{x-y-z}}.\]
(For example, if $x=30$ students average $a=60$, the top $y=5$ average $b=90$ and the bottom $z=5$ average $c=30$, the other $20$ students average $\dfrac{1800-450-150}{20}=60$.)
